\begin{table}[t]
\centering
\small
\caption{Kiểm định thống kê theo cặp giữa MADDPG và các baseline trên sum-rate (locked test bank, $N=32$; trung bình mức hạt giống, $n=8$ cặp). Pipeline Dynamic-MDP dùng paired t-test $+$ sign-flip permutation với hiệu chỉnh Holm--Bonferroni cho họ bốn so sánh chính $\{$MADDPG vs TD3, TD3-Matched, DDPG, PPO$\}$. Cohen's $d$ theo cặp đo kích thước hiệu ứng.}
\label{tab:significance}
\begin{tabular}{lcccc}
\toprule
\textbf{So sánh (sum-rate)} & \textbf{$\Delta$ trung bình} & \textbf{Cohen's $d$} & \textbf{$p_{\text{Holm}}$} & \textbf{Có ý nghĩa (5\%)} \\
\midrule
MADDPG vs TD3         & $+0{,}13 \pm 0{,}24$ & $0{,}45$ & $0{,}45$ & Không \\
MADDPG vs TD3-Matched & $+0{,}11 \pm 0{,}19$ & $0{,}47$ & $0{,}45$ & Không \\
MADDPG vs DDPG        & $+0{,}58 \pm 0{,}49$ & $0{,}98$ & $0{,}081$ & Không \\
MADDPG vs PPO         & $+1{,}10 \pm 0{,}16$ & $5{,}64$ & $\mathbf{3{,}7 \times 10^{-6}}$ & \textbf{Có} \\
\bottomrule
\end{tabular}

\vspace{0.2em}
{\footnotesize Ghi chú: Sau hiệu chỉnh Holm, MADDPG chỉ vượt PPO có ý nghĩa thống kê; so với TD3, TD3-Matched và DDPG thì \emph{không} khác biệt có ý nghĩa ở mức $5\%$. Đáng chú ý, MADDPG không vượt được TD3-Matched (biến thể tác tử đơn có \emph{cùng tổng số tham số}, xem Bảng \ref{tab:complexity}), cho thấy kiến trúc phân rã đa tác tử không mang lại lợi thế thống kê ở cấu hình này. Với DDPG, hiệu số lớn ($d=0{,}98$) nhưng chưa đạt ngưỡng sau hiệu chỉnh do cỡ mẫu $n=8$ hạn chế lực kiểm định. Kết luận này nhất quán với thí nghiệm mở rộng theo $N$ (Bảng \ref{tab:scalability}): MADDPG tương đương TD3 ở mọi $N$.\par}
\hfill {\footnotesize\itshape Nguồn: Tác giả xây dựng (2026)}
\end{table}
